\section{Conclusion}\label{sec:conclusion}

Every mechanism that stops an agent from doing an irreversible thing twice rests on an answer to ``is this the same act?''. Journals answer by position, idempotency keys by client convention, content-addressed fences by a hash. All three break when an agent re-synthesises its own arguments, and asking a language model instead offers no bound on the silent error of swallowing a distinct act.

We formalised effect identity as a decision under asymmetric loss and decomposed it into a certifier, a refuter and a residual determined by the declared schema. A fence that consults an LLM judge only on that residual, and suppresses only when the maximum score over the surviving candidates exceeds a multiplicity-corrected conformal threshold calibrated on its own jurisdiction, bounds the probability of wrong suppression for the statistic it deploys---under exchangeability, marginally over calibration draws. Refining the schema shrinks the residual and the calibration pool together: decidable safety is bought with statistical budget, and we stated that trade exactly and measured it.

On a constructed corpus of 1,066 pairs the fence recovered 94.7\% of paraphrased duplicates at 1.7\% wrong suppression, against 12.1\% for canonical keys, and the safety-critical rate stayed below $\alpha$ for three judges while deduplication did not. Typed fields made an injection class impossible; free text left another largely open. These are properties of one corpus with constructed labels, and independent human validation remains to be collected. The mechanism is a bounded semantic identity and duplicate-suppression layer with fence-level at-most-once commitment under stated assumptions; external exactly-once execution still requires a provider whose idempotency key is complete for the identity of the act.
